
et la
suite exacte des @@M0@@ obtenue par tensorisation avec @@M1@@:
@@M2@@
L'image du morphisme : @@M3@@ est
@@M4@@. Par suite, on obtient la suite exacte:
@@M5@@

On fait alors son produit tensoriel par @@M6@@; comme
@@M7@@ est un @@M8@@-Module, et que
@@M9@@ est inversible, on obtient la suite exacte:
@@M10@@
On vérifie aisément que le morphisme @@M11@@ figurant dans cette suite est le morphisme canonique, de sorte qu'elle s'identifie à la suite (3.1.2). Par définition, on a donc:
@@M12@@

\textbf{3.3.} On suppose maintenant que @@M13@@ est un schéma ayant un faisceau ample. Alors il en est de même pour @@M14@@ qui est projectif sur @@M15@@, et pour @@M16@@ et @@M17@@. Par ailleurs, @@M18@@ et @@M19@@ sont des morphismes parfaits, puisque ce sont des immersions régulières; @@M20@@ est un morphisme lisse donc parfait car @@M21@@; enfin, @@M22@@ est parfait d'après 1.9. Ces morphismes étant projectifs, on peut définir les homomorphismes additifs (IV
@@M23@@

